\documentclass[11pt]{article}

\usepackage[margin=1in]{geometry}
\usepackage{amsmath,amssymb,amsthm,mathtools}
\usepackage{array}
\usepackage{enumitem}
\usepackage[hidelinks]{hyperref}
\setlength{\emergencystretch}{3em}

\newtheorem{definition}{Definition}[section]
\newtheorem{theorem}[definition]{Theorem}
\newtheorem{corollary}[definition]{Corollary}
\newtheorem{lemma}[definition]{Lemma}
\newtheorem{proposition}[definition]{Proposition}
\newtheorem{hypothesis}[definition]{Hypothesis}
\newtheorem{conjecture}[definition]{Conjecture}
\newtheorem{remark}[definition]{Remark}

\DeclareMathOperator{\Ree}{Re}
\DeclareMathOperator{\Imm}{Im}
\DeclareMathOperator{\NTZ}{NTZ}
\DeclareMathOperator{\GRH}{GRH}
\DeclareMathOperator{\RH}{RH}
\DeclareMathOperator{\argg}{arg}
\DeclareMathOperator{\Log}{Log}
\DeclareMathOperator{\logDeriv}{logDeriv}

\title{The Helix--Hilbert--Polya Bridge: A Canonical Source Proof}
\author{Samuel Lavery}
\date{Draft proof bridge, 2026-06-13}

\begin{document}
\maketitle

\begin{abstract}
This paper gives the source-side helix proof in canonical mathematical form.  The
construction starts from the prime-generated $\pi/3$ helix, the associated unitary
phasor flow, and the Gram--von Neumann self-adjoint operator.  Zero capture is read
through the resolvent trace $-L'/L$, with multiplicity given by the logarithmic
derivative residue.  A crossing is a confluence event: the fiber phase cancels, the
amplitude reaches its threshold, the sign flips, and the next spectral harmonic is
created.  The sign flip occurs at the source midpoint $\pi/6$ in the $\pi/3$ chart;
projection to the unit-circle readout preserves the midpoint, and the logarithmic
one-dimensional readout is $1/2$.  Thus the real coordinate $1/2$ is a geometric
half-unit forced by the source crossing, while the height is retained by the lifted
phase.  The proof identifies the zero-side pole, the source crossing, and the spectral
harmonic as the same event under this projection chain.
\end{abstract}

\tableofcontents

\section{Objects and Notation}

Let $\chi$ be a Dirichlet character modulo $q$ and write
\[
  L(s,\chi)
\]
for the associated Dirichlet $L$-function.  For non-principal $\chi$, define the strip zero set
\[
  \NTZ(\chi)
  :=
  \{\rho\in\mathbb C : 0<\Ree\rho<1,\; L(\rho,\chi)=0\}.
\]
For a channel $\chi$, write
\[
  \GRH(\chi)
  :=
  \forall \rho\in\NTZ(\chi),\quad \Ree\rho=\frac12 .
\]
The all non-principal channel closure statement is
\[
  \GRH_{\ne 1}
  :=
  \forall q,\;\forall\chi \pmod q,\quad \chi\ne1\Longrightarrow\GRH(\chi).
\]
The principal zeta/L1 channel is treated separately through eta regularization.
For zeta, let
\[
  \NTZ_\zeta
  :=
  \{\rho\in\mathbb C : 0<\Ree\rho<1,\; \zeta(\rho)=0\}.
\]

\section{State of the Art and Contribution}

The classical analytic theory begins with Riemann's use of $\zeta(s)$, its
functional equation, and its zero set in the explicit study of primes
\cite{Riemann1859,Edwards1974,Titchmarsh1986}.  For Dirichlet and related
$L$-functions, the standard arithmetic input is the Euler product, its
logarithmic derivative, and the von Mangoldt prime-power trace
\cite{Davenport2000,MontgomeryVaughan2007,IwaniecKowalski2004}.  In this
paper those facts are not used as a strip-convergent Euler product.  They are
used as the source discriminator: the fiber being projected must be the
prime-generated Euler-product fiber, not merely a function with a matching
functional equation.

The spectral viewpoint is usually formulated as a search for an operator or
trace formula whose spectral side reads the critical zeros.  Representative
models include the $H=xp$ program, noncommutative trace formulas, and random
matrix models for zero statistics and central values
\cite{BerryKeating1999,Connes1999,KeatingSnaith2000}.  The present construction
uses a different source object.  The operator side is the Gram--von Neumann
normal form $G_\infty=B_\infty^*B_\infty$ on the source Hilbert space; the
geometric side is the prime-generated $\pi/3$ helix with no radial drift.

The Davenport--Heilbronn phenomenon shows why a functional-equation-level
comparison is insufficient for this proof strategy \cite{DavenportHeilbronn1936}.
The readout chain therefore includes an Euler-product discriminator before the
Pythagorean midpoint step.  This is the point at which an object lacking the
prime-generated source cannot enter the same projection argument.

The contribution of this paper is the following canonical bridge:
\begin{enumerate}[label=(C\arabic*)]
\item construct the placed source fiber directly from the $\pi/3$ helix and
the prime-generated channel data;
\item identify the closed-form amplitude accumulation $C^{-s}L(s,\chi)$ and
the resolvent trace $C^{-s}(-L'/L)(s,\chi)$ as two readouts of the same source;
\item model zero capture as a crossing confluence where phase cancellation,
amplitude threshold, sign flip, trace residue, and new harmonic creation occur
simultaneously;
\item prove that the source sign flip forces the midpoint $\pi/6$ in the
$\pi/3$ chart, and that the logarithmic one-dimensional readout is
$(3/\pi)(\pi/6)=1/2$.
\end{enumerate}

\section{Proof Roadmap}

The proof has four layers.

\begin{enumerate}[label=(P\arabic*)]
\item \textbf{Source construction.}
The integer fibers are placed on the $\pi/3$ helix by the area law
$R_n^2=Cn$, with $C=(\pi/3)^2$.  Unique factorization supplies the prime
coordinates, the channel weights, and the Euler-product discriminator.

\item \textbf{Trace and accumulation.}
The placed fiber is the amplitude accumulation $C^{-s}L(s,\chi)$, while the
resolvent trace is the prime-power accumulation
$C^{-s}(-L'/L)(s,\chi)$.  A captured zero is therefore read as a pole of the
resolvent trace, with residue equal to its multiplicity up to the nonzero
geometric gauge.

\item \textbf{Crossing geometry.}
At a source crossing the channel phase cancels, the amplitude reaches a
threshold, the sign flips, and the next harmonic begins.  The sign flip occurs
on the source midpoint axis, which is $\pi/6$ in the $\pi/3$ chart.

\item \textbf{Projection readout.}
The three-dimensional crossing projects to the unit circle without changing the
midpoint angle.  The logarithmic one-dimensional readout sends
$\pi/6$ to $(3/\pi)(\pi/6)=1/2$ while the lifted phase retains the height.
\end{enumerate}

The subsequent sections supply these layers in the order needed to reproduce
the argument: source geometry, Euler-product generation, trace normal forms,
standing readout, crossing costs, the R1--R12 infrastructure, the crossing
definition, and finally the canonical closure.

\section{\texorpdfstring{The $\pi/3$ Helix Geometry}{The pi/3 Helix Geometry}}

\subsection{Canonical Scale}

We use the canonical $\pi/3$ chart
\[
  A=\frac{\pi}{6},\qquad ds=\frac{\pi}{3},\qquad
  C=2A\,ds=\left(\frac{\pi}{3}\right)^2.
\]
The exact area-law placement is
\[
  R_n^2=Cn,\qquad R_n=\frac{\pi}{3}\sqrt n .
\]
At the source level the data are the $\pi/3$ coordinates, the channel weights
coming from residue counting, and the harmonic phase on the helix.  The fiber is
derived from that source.  No one-dimensional continuation formula is used to
create it.  The standard symbols $L$, $\zeta$, and Hurwitz tails below occur
because parts of the derived fiber identify with those representations after the
geometric fiber has already been placed.
Thus the three-dimensional source fiber atom at height $t$, for
$s=\frac12+it$, is
\[
  a_n(s)
  =
  w(n)\frac{1}{R_n}\exp\{-it\Log(R_n^2)\}
  =
  w(n)(Cn)^{-s}.
\]

\begin{lemma}[Area-law closed form]
For any channel weights $w(n)$ for which the corresponding Dirichlet series is defined,
\[
  F_w(s)
  :=
  \sum_{n\ge1} w(n)(Cn)^{-s}
  =
  C^{-s}\sum_{n\ge1} w(n)n^{-s}.
\]
\end{lemma}

\begin{proof}
Use $R_n^2=Cn$ and $s=\frac12+it$:
\[
  R_n^{-1}e^{-it\Log(R_n^2)}
  =
  (R_n^2)^{-1/2}(R_n^2)^{-it}
  =
  (R_n^2)^{-s}
  =
  (Cn)^{-s}.
\]
Summing gives the result.
\end{proof}

\subsection{Non-principal Character Channels}

For a non-principal Dirichlet character $\chi$ modulo $q$, the channel weights are
\[
  w(n)=\chi(n).
\]
Therefore the source channel sum is
\[
  F_\chi(s)
  =
  \sum_{n\ge1}\chi(n)(Cn)^{-s}.
\]
In the projected one-dimensional notation this same harmonic sum is
\[
  F_\chi(s)=C^{-s}L(s,\chi).
\]
Residue-wise, with zero residues omitted, the same placed channel may be
reindexed as
\[
  F_\chi(s)
  =
  C^{-s}q^{-s}\sum_{r\bmod q}\chi(r)\,
  \zeta\!\left(s,\frac rq\right).
\]
For $\chi_3$,
\[
  \chi_3(1)=1,\qquad \chi_3(2)=-1,\qquad \chi_3(0)=0,
\]
and hence
\[
  F_{\chi_3}(s)
  =
  C^{-s}3^{-s}
  \left[
    \zeta\!\left(s,\frac13\right)
    -
    \zeta\!\left(s,\frac23\right)
  \right].
\]

\subsection{The Principal/Zeta Channel}

For the zeta/L1 channel the fiber is eta-regularized:
\[
  w(n)=(-1)^{n+1}.
\]
The source channel sum is
\[
  F_{L1}(s)=\sum_{n\ge1}(-1)^{n+1}(Cn)^{-s}.
\]
In the projected one-dimensional notation the eta factor is
\[
  \eta(s)
  =
  \sum_{n\ge1}(-1)^{n+1}n^{-s}
  =
  (1-2^{1-s})\zeta(s).
\]
Thus the L1 fiber is
\[
  F_{L1}(s)
  =
  C^{-s}\eta(s)
  =
  C^{-s}(1-2^{1-s})\zeta(s).
\]
Equivalently, in residue form modulo $2$,
\[
  F_{L1}(s)
  =
  C^{-s}2^{-s}
  \left[
    \zeta\!\left(s,\frac12\right)-\zeta(s,1)
  \right].
\]
On the critical line $\Ree s=\frac12$, the eta factor is nonzero:
\[
  1-2^{1-s}\ne0,
\]
because $2^{1-s}=1$ would imply $2^{1-\Ree s}=1$, hence $\Ree s=1$.

\subsection{FTA and Euler Product Generation}

The source helix is prime-generated.  Its integer fibers are generated from
prime coordinates by unique factorization, and the Euler product is the
convergence-half-plane readout of that multiplicative geometry
\cite{Davenport2000,MontgomeryVaughan2007}.

\begin{proposition}[FTA helix character]
Fix a prime-angle assignment $\theta:\mathbb N\to\mathbb R$ and define the FTA
winding angle
\[
  \Theta_\theta(n)
  =
  \sum_{p^e\parallel n} e\,\theta(p).
\]
For $m,n\ge1$,
\[
  \Theta_\theta(mn)
  =
  \Theta_\theta(m)+\Theta_\theta(n).
\]
Consequently the unit-circle winding
\[
  \omega_\theta(n)=\exp(i\Theta_\theta(n)),
\]
is multiplicative:
\[
  \omega_\theta(mn)=\omega_\theta(m)\omega_\theta(n).
\]
For a Dirichlet channel $\chi$, the channel helix character
\[
  \Psi_{\chi,\theta}(n)=\chi(n)\omega_\theta(n)
\]
is multiplicative:
\[
  \Psi_{\chi,\theta}(mn)
  =
  \Psi_{\chi,\theta}(m)\Psi_{\chi,\theta}(n)
  \qquad(m,n\ge1).
\]
\end{proposition}

\begin{proof}
By unique factorization,
\[
  v_p(mn)=v_p(m)+v_p(n)
\]
for every prime $p$.  Summing $e\,\theta(p)$ over the finite factorization
support gives the displayed additivity of $\Theta_\theta$.  Exponentiating
converts addition of angles into multiplication on the unit circle.  Since
Dirichlet characters are multiplicative, the product
$\chi(n)\omega_\theta(n)$ is multiplicative as well.

\end{proof}

\begin{proposition}[Log bridge and Euler product readout]
The logarithm enters only at the external analytic bridge.  With bridge angle
\[
  \theta_\gamma(p)=\gamma\log p,
\]
the FTA identity
\[
  \log n=\sum_{p^e\parallel n} e\log p
\]
gives
\[
  \Theta_{\theta_\gamma}(n)=\gamma\log n,
  \qquad
  \omega_{\theta_\gamma}(n)=n^{i\gamma}.
\]
In the convergence half-plane $\Ree s>1$, the same multiplicativity is read as
an Euler product:
\[
  \operatorname{HelixSource}_\chi^C(s)
  =
  \sum_{n\ge1}\chi(n)(Cn)^{-s}
  =
  C^{-s}\prod_p\left(1-\chi(p)p^{-s}\right)^{-1}.
\]
Taking the logarithmic derivative gives the prime-power trace generated by the
Euler product:
\[
  \mathcal T_\chi^C(s)
  =
  C^{-s}\sum_{n\ge1}\Lambda(n)\chi(n)n^{-s}
  =
  C^{-s}\left(-\frac{L'}{L}(s,\chi)\right).
\]
The winded prime-field bridge is
\[
  \sum_{n\ge1}\chi(n)\Lambda(n)\omega_{\theta_\gamma}(n)n^{-s}
  =
  -\frac{L'}{L}(s-i\gamma,\chi),
  \qquad \Ree s>1.
\]
Where the product representation converges, it also gives the Euler-product
edge fact
\[
  \Ree s\ge1\Longrightarrow L(s,\chi)\ne0
  \qquad(\chi\ne1).
\]
\end{proposition}

\begin{proof}
The identity
\[
  \log n=\sum_{p^e\parallel n} e\log p
\]
is the logarithmic image of unique factorization.  Substituting
$\theta_\gamma(p)=\gamma\log p$ into the FTA angle gives
$\Theta_{\theta_\gamma}(n)=\gamma\log n$, hence
$\omega_{\theta_\gamma}(n)=n^{i\gamma}$.  In $\Ree s>1$, absolute convergence
allows the completely multiplicative integer series to factor into its prime
Euler product.  Applying the logarithmic derivative to the Euler product gives
the von Mangoldt prime-power trace.  Multiplication by
$\omega_{\theta_\gamma}(n)=n^{i\gamma}$ shifts the trace from $s$ to
$s-i\gamma$.

\end{proof}

Thus composites are not independent sources; they are generated by their prime
coordinates through FTA.  This is the structural distinction from a
functional-equation-only construction.  The product convergence is
half-plane-limited; the FTA multiplicativity of the helix character itself is
not height-limited.

For the principal L1 channel, the same prime generation is present with
$\chi=1$.  The working fiber is eta-regularized,
\[
  F_{L1}(s)=C^{-s}(1-2^{1-s})\zeta(s),
\]
and the zero readout uses the gauged zeta resolvent trace
\[
  C^{-s}\left(-\frac{\zeta'}{\zeta}(s)\right).
\]
The eta factor handles the principal-channel pole; it does not replace prime
generation.

\subsection{Euler-Product Discriminator}

The previous propositions supply a discriminator for the source object rather
than a standalone zero-location argument.
The logarithmic Euler product
\[
  \log\zeta(s)
  =
  \sum_p\sum_{k\ge1}\frac{1}{k}p^{-ks}
\]
is a convergence-half-plane identity.  It is not evaluated at a strip zero.
The role of the Euler product in the closure is instead to certify that the
source fiber being projected is the prime-generated Dirichlet/zeta fiber, not
only an object with a matching functional equation.

\begin{definition}[Euler-product discriminator]
For a non-principal Dirichlet channel $\chi$, the Euler-product discriminator
\[
  \mathcal E_\chi
\]
is the following package:
\begin{enumerate}[label=(E\arabic*)]
\item the source fiber factors into the prime Euler product in $\Ree s>1$:
\[
  \operatorname{HelixSource}_\chi^C(s)
  =
  C^{-s}\prod_p(1-\chi(p)p^{-s})^{-1};
\]
\item the winded prime-power trace is the shifted logarithmic derivative:
\[
  \sum_{n\ge1}\chi(n)\Lambda(n)\omega_{\theta_\gamma}(n)n^{-s}
  =
  -\frac{L'}{L}(s-i\gamma,\chi),
  \qquad \Ree s>1;
\]
\item the Euler-product edge gives nonvanishing on $\Ree s\ge1$:
\[
  \Ree s\ge1\Longrightarrow L(s,\chi)\ne0 .
\]
\end{enumerate}
\end{definition}

\begin{proposition}[Euler product in the readout chain]
The Pythagorean readout consumes the Euler-product discriminator before it
converts a projected source event into the midpoint statement.  In mathematical
terms, the readout step uses the package
\[
  \mathcal E_\chi,\qquad
  \operatorname{FiberCaptureEvent}_\chi(\rho),\qquad
  \operatorname{SourceFiberEvent}_\chi(\rho),\qquad
  \text{no radial drift},\qquad
  \operatorname{Proj}_\chi(\rho)
\]
to obtain
\[
  \|w(\rho)\|^2=1.
\]
The Pythagorean unit-circle identity then gives the midpoint conclusion
\[
  \Ree\rho=\frac12.
\]
\end{proposition}

\begin{proof}
The readout chain first builds $\mathcal E_\chi$ from the Dirichlet Euler
product and winded von Mangoldt trace.  Only then does the projection chain
provide the unit-circle/Pythagorean readout.  Therefore a
functional-equation-only model cannot pass through the same readout step unless
it supplies the same discriminator package.  This is the falsifier check: if
the source object is replaced by a Davenport--Heilbronn-type readout with the
same symmetry but no Euler-product source package, the theorem is not
applicable at the readout step.
\end{proof}

\section{Source Accumulation and Resolvent Trace}

The placed fiber accumulation and the resolvent-trace accumulation are distinct
but linked closed forms.  This section records both normal forms before the
standing readout and crossing geometry are introduced.

\subsection{Polya Operator Normal Form}

For a non-principal channel, the placed source fiber is
\[
  F_\chi^C(s)
  =
  \sum_{n\ge1}\chi(n)(Cn)^{-s}
  =
  C^{-s}L(s,\chi).
\]
This is the channel amplitude accumulation: it is the harmonic sum of the
source fiber terms in the $\pi/3$ area chart.

\begin{lemma}[No radial drift in the accumulation normal form]
The helix source has no radial drift by construction.  For each source atom
\[
  a_n(t)=w(n)R_n^{-1}\exp\{-it\Log(R_n^2)\},
\]
the radial size is independent of $t$:
\[
  |a_n(t)|
  =
  \frac{|w(n)|}{R_n}
  =
  \frac{|w(n)|}{(\pi/3)\sqrt n}
  \qquad (C=(\pi/3)^2).
\]
\end{lemma}

\begin{proof}
For real $t$ and positive $R_n^2$,
\[
  \left|\exp\{-it\Log(R_n^2)\}\right|=1.
\]
The helix flow therefore changes only phase.  The radial coordinate is fixed by
the source placement $R_n^2=Cn$, so the accumulated source is no-drift in the
radial direction before projection.
\end{proof}

This no-drift normal form is the invariant used by the Polya induction.  Moving
from one crossing to the next advances the phase ledger and harmonic threshold,
not the radial placement.  The same $\pi/3$ source geometry therefore persists
through the infinite ladder of crossings.  In the R1--R12 summary below, this
is the source-side input for the geometric crossing correspondence and the
resolvent-trace identification.

\begin{definition}[Admissible source harmonic domain]
For $s=\sigma+it$ and a finite set $F\subset\mathbf N$, the admissible finite
source harmonic is
\[
  H_{F,\chi}(s)
  =
  \sum_{n\in F}\chi(n)n^{-\sigma}U(t,n)^{-1},
  \qquad
  U(t,n)=\exp(it\log n).
\]
For the initial segment $F=\{0,\ldots,M\}$ this is the finite source chain.
The admissible source domain is the increasing ladder obtained from these
finite source harmonics by threshold induction, with the $\pi/3$ coordinate
chart applied at every crossing.  No Green--Helmholtz energy norm is part of
this source domain.
\end{definition}

\begin{proposition}[Source operator normal form]
The Polya source flow is the diagonal unitary flow
\[
  U(t,n)=\exp(it\log n),
  \qquad
  \|U(t,n)\|=1,
\]
with real generator
\[
  H(n)=\log n,\qquad \Im H(n)=0.
\]
Consequently a conserved nonzero source amplitude has zero radial drift.
\end{proposition}

\begin{proof}
The flow is unitary because for real $t$,
\[
  |e^{it\log n}|=1
\]
for every $n$.  Its Stone generator is the real diagonal multiplication
operator
\[
  H(n)=\log n.
\]
Since the generator has no imaginary part and the flow is unitary, a conserved
nonzero amplitude cannot acquire an exponential radial factor.  The no-drift
lemma then gives $\Ree\lambda=0$ for the source drift rate.
\end{proof}

\begin{theorem}[The helix is its own continuation into the strip]
For a non-principal channel $\chi$, define the finite source phasor chain
\[
  S_M(s,\chi)
  =
  \sum_{0\le n\le M}
    \chi(n)n^{-\sigma}U(t,n)^{-1},
  \qquad
  s=\sigma+it,\quad U(t,n)=\exp(it\log n).
\]
This is the same three-dimensional helix fiber, projected through the
area-law/log bridge.  For every $s$ with $\sigma>0$ there is a constant
$C_{\chi,s}>0$ such that
\[
  \left|S_M(s,\chi)-L(s,\chi)\right|
  \le
  C_{\chi,s}M^{-\sigma}
  \qquad(M\ge1).
\]
Thus the helix source chain continues itself from the convergent half-plane
into the critical strip.  At a zero $\rho$ with $0<\Ree\rho<1$,
\[
  |S_M(\rho,\chi)|
  \le
  C_{\chi,\rho}M^{-\Ree\rho},
\]
so the zero is read as a source fiber cancellation event.  Moreover the
rescaled chain exposes the height-free source ledger:
\[
  \left|
    S_M(\rho,\chi)M^\rho
    -
    \left(A_\chi(M)-L(0,\chi)\right)
  \right|
  \le
  C'_{\chi,\rho}M^{-1},
  \qquad
  A_\chi(M)=\sum_{0\le n\le M}\chi(n).
\]
\end{theorem}

\begin{proof}
The phasor-flow form of the source chain is the displayed finite sum.  The
continuation estimate is obtained by summation by parts using the bounded
partial sums of a non-principal character.  At a zero, substituting
$L(\rho,\chi)=0$ gives the cancellation estimate.  Multiplying the estimate by
$M^\rho$ gives the height-free ledger displayed above.
This is not the divergent Euler product in the strip; it is the same harmonic
source chain continued by the closure ledger.  The zeta/L1 channel uses the
parallel eta-regularized fiber and standing readout.
\end{proof}

\subsection{Projection and Threshold Normal Forms}

\begin{proposition}[Pythagorean midpoint readout]
For a projected source fiber event, the Pythagorean unit-circle readout forces
\[
  \operatorname{normSq}(w(\rho))=1
  \Longleftrightarrow
  \Ree\rho=\frac12.
\]
At the $\pi/3$ source scale this is the same midpoint statement
\[
  \operatorname{arcChart}(\Ree\rho)=\frac{\pi}{6}.
\]
Equivalently, if $z_2(\rho)\in S^1$ is the two-dimensional unit-circle
projection of the crossing, written in the $\pi/3$ angular chart, then the
one-dimensional readout is the lifted logarithm
\[
  \ell(\rho)
  =
  \frac{3}{\pi}\frac{1}{i}\Log_{\mathcal H} z_2(\rho),
  \qquad
  z_2(\rho)=\exp\!\left(i\,\frac{\pi}{3}\ell(\rho)\right),
\]
where $\Log_{\mathcal H}$ is the continuous helix branch selected by the
crossing ledger.  Thus the unit-circle projection is still a $\pi/3$-coordinate
object, and the 1D value is obtained by taking its logarithm and applying the
same scale conversion.
\end{proposition}

\begin{proof}
For $w(\rho)=1-1/\rho$,
\[
  |w(\rho)|^2=1
  \Longleftrightarrow
  \left(1-\frac1\rho\right)\left(1-\frac1{\overline\rho}\right)=1
  \Longleftrightarrow
  \rho+\overline\rho=1
  \Longleftrightarrow
  \Ree\rho=\frac12.
\]
The $\pi/3$ coordinate conversion is
\[
  \operatorname{arcChart}(x)=\frac{\pi}{3}x,
  \qquad
  \operatorname{arcChart}^{-1}(u)=\frac3\pi u.
\]
The value $1/2$ is therefore the unit-readout midpoint forced by the source
sign-change/no-drift crossing after source projection, not an independent
condition on the zero set.
\end{proof}

\begin{proposition}[Threshold ladder]
The first source crossing has amplitude cost $\pi/2$, and each subsequent
harmonic threshold advances by $\pi$:
\[
  A_0^+=\frac{\pi}{2},
  \qquad
  A_{k+1}^+=A_k^++\pi .
\]
At a threshold $E(t)=n\pi$, the harmonic count is exactly $n$.
\end{proposition}

\begin{proof}
The positive ladder is $A_k^+=(k+\frac12)\pi$, so
$A_0^+=\pi/2$ and $A_{k+1}^+-A_k^+=\pi$.  If the harmonic count is defined by
$\lfloor E(t)/\pi\rfloor$ at threshold, then $E(t)=n\pi$ gives count $n$.
The channel cancellation law is the equality of the positive and negative
parts of the signed channel sum.
\end{proof}

\begin{remark}[What does and does not enter the Polya chain]
The operator used here is the Gram/von Neumann self-adjoint operator
$G_\infty=B_\infty^*B_\infty$.  Green--Helmholtz energy and finite Dirac
operator models are separate analytic models and are not part of this proof
chain.  The Polya chain used here is exactly the Gram/von Neumann
self-adjoint operator, phasor-flow unitarity, real generator/no-drift,
threshold ladder, source cancellation, resolvent pole, $\pi/3$ chart, and
Pythagorean readout path above.
\end{remark}

\begin{proposition}[Exact multiplicity in the trace]
Assume the boundary trace identity
\[
  T(z)=
  -\frac{L'}{L}\!\left(\frac12+iz,\chi\right).
\]
If $\rho\in\NTZ(\chi)$ has order $m$, then at
$z=\operatorname{poleParam}(\rho)$ the resolvent trace has a simple principal
part with residue $i\,m$.  In the $s$-variable and $\pi/3$ gauge this is the
same residue law
\[
  \operatorname*{Res}_{s=\rho}\mathcal T_\chi^C(s)=-mC^{-\rho}.
\]
\end{proposition}

\begin{proof}
If $L(s,\chi)=(s-\rho)^m g(s)$ with $g(\rho)\ne0$, then
$L'/L=m/(s-\rho)+g'/g$.  Under the boundary coordinate
$s=\frac12+iz$, the pole parameter changes the residue by the coordinate
Jacobian, giving the displayed $i\,m$ in the $z$-trace.  The multiplicity is
therefore read from
the Laurent residue of the resolvent trace, not from a numerical fit.
\end{proof}

\begin{proposition}[Source normal form for crossing correspondence]
The source normal form used by the crossing correspondence has three mathematical
components.  First, radial drift is zero under source conservation.  Second,
the purchase-height ladder satisfies
\[
  E(\operatorname{purchaseHeight}(n))=n\pi,\qquad
  \operatorname{harmonicCount}(E,\operatorname{purchaseHeight}(n))=n,
\]
with rigidity for every sequence satisfying the same conversion law.  Third,
each purchase crossing is identified with the corresponding midpoint geometric
zero and unit-circle readout.
\end{proposition}

\begin{proof}
The no-radial-drift normal form fixes the source placement $R_n^2=Cn$ at every
harmonic threshold.  Threshold existence and the quantum ladder locate each
crossing.  Ladder exhaustion and rigidity show that the purchase sequence is
the only possible crossing sequence.  The $\pi/3$ chart supplies the midpoint
coordinate, and the resolvent-pole side is supplied by the
logarithmic-derivative pole law.
\end{proof}

\subsection{Trace Residues and Completed Readout}

The resolvent trace is the logarithmic, prime-power-weighted accumulation.  In
the convergent half-plane,
\[
\begin{aligned}
  \mathcal T_\chi^C(s)
  &:=
  \sum_{n\ge1}\Lambda(n)\chi(n)(Cn)^{-s}  \\
  &=
  C^{-s}\sum_{n\ge1}\Lambda(n)\chi(n)n^{-s} \\
  &=
  C^{-s}\left(-\frac{L'}{L}(s,\chi)\right).
\end{aligned}
\]
Here $\Lambda$ is the von Mangoldt weight.  Thus the same source geometry has
two closed-form accumulations: the fiber sum $C^{-s}L(s,\chi)$ and the
resolvent trace $C^{-s}(-L'/L)(s,\chi)$.
At the canonical scale $C=(\pi/3)^2$, the continued trace is
\[
  \mathcal T_\chi^{\pi/3}(s)
  =
  \left(\left(\frac{\pi}{3}\right)^2\right)^{-s}
  \left(-\frac{L'}{L}(s,\chi)\right),
\]
which is the canonical $\pi/3$ continued trace.

The literal logarithmic derivative of the placed fiber differs only by a
regular gauge term:
\[
  -\frac{d}{ds}\log F_\chi^C(s)
  =
  \log C-\frac{L'}{L}(s,\chi).
\]
The additional term $\log C$ has no pole.  The multiplicative geometric gauge
$C^{-s}$ is never zero for $C>0$.  Therefore neither gauge changes the crossing
locations or creates/removes trace poles; the multiplicative gauge only rescales
the trace residue by a nonzero geometric factor.

This gives the local residue law at a capture event.  If $\rho$ is a zero of
$L(s,\chi)$ of order $m$ and
\[
  L(s,\chi)=(s-\rho)^m g(s),\qquad g(\rho)\ne0,
\]
then near $\rho$
\[
  -\frac{L'}{L}(s,\chi)
  =
  -\frac{m}{s-\rho}-\frac{g'}{g}(s),
\]
and hence
\[
  \mathcal T_\chi^C(s)
  =
  -\frac{mC^{-\rho}}{s-\rho}
  +\text{a term finite at }\rho .
\]
So the trace pole has residue $-mC^{-\rho}$.  The integer $m$ is the spectral
multiplicity; the factor $C^{-\rho}$ is the nonzero $\pi/3$ geometric scaling.
At a crossing, this is the residue consumed by the resolvent trace while the
source fiber amplitude crests, the channel phase cancels, and the sign changes.
The Polya threshold condition is therefore the simultaneous event
\[
\begin{gathered}
  F_\chi^C(\rho)=0,\qquad
  \operatorname*{Res}_{s=\rho}\mathcal T_\chi^C(s)=-mC^{-\rho},\\
  E_\chi(\Imm\rho)=A_k^\pm,
  \qquad
  A_0^\pm=\pm\frac{\pi}{2},\qquad
  A_{k+1}^\pm-A_k^\pm=\pm\pi .
\end{gathered}
\]
Here $E_\chi$ is the source accumulation ledger used by the Polya chain.  The
jump from one level $A_k^\pm$ to the next is the threshold that starts the next
harmonic on the spectral readout.  Thus the closed-form resolvent residue, the
crossing sign change, and the creation of the next harmonic are the same source
event viewed in three ledgers: trace, standing wave, and harmonic ladder.

For primitive non-principal $\chi$, the completed zero-side resolvent trace is
the partial-fraction accumulation
\[
  \mathcal R_\chi(s)
  =
  \sum_{\rho\in\NTZ(\chi)}
    m_\rho
    \left(\frac{1}{s-\rho}+\frac{1}{\rho}\right),
\]
where $m_\rho$ is the vanishing order of the completed $L$-function at $\rho$.
The closed form is
\[
  \logDeriv \Lambda_\chi(s)
  =
  A+\mathcal R_\chi(s)
  \qquad (s\notin\NTZ(\chi)),
\]
for a constant $A$.  Splitting the completed function
$\Lambda_\chi=\gamma_\chi L(s,\chi)$ gives, off the zeros in the strip,
\[
  -\frac{L'}{L}(s,\chi)
  =
  -A-\mathcal R_\chi(s)+\logDeriv \gamma_\chi(s).
\]
Thus the prime-power accumulation, the zero-side resolvent trace, and the
continued source trace are the same readout up to the constant and
Archimedean gauge terms, which are regular away from their own gauge poles.

For the L1/zeta channel,
\[
  F_{L1}^C(s)=C^{-s}(1-2^{1-s})\zeta(s).
\]
Off $\Ree s=1$, the eta factor is nonzero, so a zero of the L1 fiber is a zero
of $\zeta$.  The corresponding continued helix resolvent trace is
\[
  \mathcal T_{L1}^C(s)
  =
  C^{-s}\left(-\frac{\zeta'}{\zeta}(s)\right).
\]
The eta factor regularizes the principal fiber; it is not inserted into the
zero-pole readout, because the pole is the logarithmic derivative of the zeta
factor whose zeros are being read.  The raw trace is
\[
  -\frac{\zeta'}{\zeta}(s),
\]
and the $\pi/3$ zeta fiber zero has both a raw and a gauged resolvent-trace
pole, since multiplication by the nonzero gauge $C^{-s}$ cannot remove a pole.

\section{Standing Readout and Projection to One Dimension}

This is the layer where the placed three-dimensional fiber is read in the
lower-dimensional coordinate chart.  The functions below are readouts of the
source geometry and channel counts, with the coordinate phase correction kept
explicit.

For $s=\frac12+it$, the projected complex fiber is
\[
  F_\chi(t)=F_\chi\!\left(\frac12+it\right).
\]
The real standing readout has the form
\[
  Z_\chi(t)
  =
  \Ree\left(
    e^{i\Theta_\chi(t)}
    F_\chi\!\left(\frac12+it\right)
  \right),
\]
where
\[
  \Theta_\chi(t)
  =
  \argg\Gamma\!\left(\frac{1/2+a}{2}+\frac{it}{2}\right)
  +\frac t2\log\frac q\pi
  +t\log C
  +\alpha_\chi .
\]
Here $a\in\{0,1\}$ is the parity and $\alpha_\chi$ is the character/root-number gauge.
For the L1/zeta eta-regularized channel one subtracts the eta gauge
\[
  \argg(1-2^{1-s}).
\]

The term $t\log C$ is not a new oscillation.  It is the coordinate correction from
\[
  \Log(R_n^2)=\Log(Cn)=\log n+\log C.
\]
Omitting it rotates the real readout into the wrong coordinate chart.

\subsection{Phase Conventions}

All logarithms of positive real quantities are real logarithms.  Thus
\[
  \Log(R_n^2)=\log(Cn),\qquad
  \log q,\qquad \log C
\]
use the real branch.  The gamma phase in the standing readout is the continuous
phase
\[
  \vartheta_{\chi}(t)
  =
  \Imm\log\Gamma\!\left(\frac{1/2+a+it}{2}\right)
  +\frac t2\log\frac q\pi ,
\]
with the branch chosen continuously along the vertical line
$z(t)=(1/2+a+it)/2$.  The root-number phase is fixed by
\[
  \tau(\chi)=\sum_{r=1}^{q}\chi(r)e^{2\pi i r/q},\qquad
  \varepsilon_\chi=\frac{\tau(\chi)}{i^a\sqrt q},\qquad
  \alpha_\chi=-\frac12\argg(\varepsilon_\chi).
\]
For the principal zeta/L1 channel, $\alpha_\chi=0$ and the two-channel
eta readout uses the continuous phase
\[
  \phi_\eta(t)
  =
  \argg\!\left(1-2^{1-(1/2+it)}\right),
\]
so that
\[
  \Theta_{L1}(t)
  =
  \vartheta_\zeta(t)+t\log C-\phi_\eta(t).
\]
For non-principal $\chi$,
\[
  \Theta_\chi(t)=\vartheta_\chi(t)+t\log C+\alpha_\chi.
\]

\section{Harmonic Cost and the Quantum Ladder}

The full source geometry is double ended.  The midpoint is the origin of the
fiber, and the reflected side is an anti-helix carrying the opposite orientation.
The proof uses the positive oriented side of the helix for simplicity; the
reflected anti-helix is not used in the proof below.  It is useful for
computational verification: the reflected side supplies the symmetric
pre-origin geometry needed to build a high-precision tail before the counting
ledger starts at the origin.  Its cost law is the same with the sign reversed.

The crossing levels in the projected unit-circle readout are therefore
\[
  A_k=\left(k+\frac12\right)\pi,\qquad k\in\mathbb Z.
\]
Equivalently, for $m\ge0$ the two oriented ladders are
\[
  A_m^+=\left(m+\frac12\right)\pi,\qquad
  A_m^-=-\left(m+\frac12\right)\pi .
\]
On the positive helix side used in the proof,
\[
  A_0=\frac\pi2,\qquad A_{k+1}=A_k+\pi.
\]
Thus the first crossing is the only crossing on that side whose cost from the
origin is a half quantum:
\[
  |A_0^+|=\frac\pi2.
\]
Every additional crossing on that same side costs one full quantum:
\[
  A_{m+1}^+-A_m^+=\pi\qquad(m\ge0).
\]
The reflected anti-helix obeys the same rule:
\[
  |A_0^-|=\frac\pi2,\qquad |A_{m+1}^- - A_m^-|=\pi\qquad(m\ge0).
\]
So the half-quantum occurs once at the first crossing on each oriented side of
the double-ended helix; all subsequent crossings consume full $\pi$ increments.

\begin{conjecture}[Origin/Tate interface]
If this helix model has a Tate-thesis-style global/local interpretation, the
natural place for that connection is the shared origin of the double-ended
helix.  The origin is where the two reflected orientations meet before the
one-dimensional readout separates them into positive and negative crossing
ladders \cite{Tate1967}.
\end{conjecture}

\section{Hilbert--Polya Source Infrastructure: R1--R12}

This section lists the mathematical infrastructure used by the proof.  R1--R11
are the operator, flow, source, zero, crossing, coordinate, and spectral-reality
facts.  R12 is the counting and midpoint-identification package used to close
the projection readout.

\begin{enumerate}[label=R\arabic*., leftmargin=2cm]
\item \textbf{Hilbert space.}
The source space is
\[
  \ell^2(\mathbb N;\mathbb C).
\]
\item \textbf{Self-adjoint Gram operator.}
The Gram operator is the von Neumann normal form
\[
  G_\infty=B_\infty^*B_\infty .
\]
For any closed densely defined operator $T$ on a Hilbert space, the theorem
\[
  T^*T \text{ is self-adjoint}
\]
holds \cite{ReedSimon1980}.  The helix Gram operator is this concrete
instantiation.  Eigenvalues of
the symmetric/self-adjoint operator are real in R11.

\item \textbf{Earned unitary phasor flow.}
The FTA/Euler phasor flow is
\[
  U(t)(n)=n^{it},\qquad |U(t)(n)|=1,
\]
with
\[
  U(s+t)(n)=U(s)(n)U(t)(n).
\]
\item \textbf{Real generator and no-drift principle.}
The generator is
\[
  H(n)=\log n\in\mathbb R,\qquad U(t)(n)=e^{itH(n)}.
\]
The source no-drift lemma is:
if
\[
  e^{\Ree\lambda\,\tau}c=c\quad\text{for all }\tau,\qquad c\ne0,
\]
then
\[
  \Ree\lambda=0.
\]
\item \textbf{Analytic $L$-function object.}
For every character channel, the associated $L$-function is differentiable away from
the point $s=1$:
\[
  s\ne1\Longrightarrow
  \operatorname{DifferentiableAt}_{\mathbb C}L(s,\chi).
\]
For non-principal channels this is the Dirichlet closure ledger.  For the
principal L1/zeta channel the eta-regularized fiber carries the same zero
readout:
\[
  F_{L1}(s)=C^{-s}(1-2^{1-s})\zeta(s),
\]
with the eta factor nonzero off $\Ree s=1$.

\item \textbf{Zeros induce fiber capture events.}
At a zero $L(s,\chi)=0$ with $0<\Ree s$,
\[
  \|S_M(s,\chi)\|\le C_s M^{-\Ree s}.
\]
At every $\rho\in\NTZ(\chi)$,
\[
  -\frac{L'}{L}(s,\chi)
\]
has no finite limit at $\rho$ in the punctured neighborhood.
This is the R6 resonance clause:
\[
  \rho\in\NTZ(\chi)\Longrightarrow
  -L'/L\text{ has a pole at }\rho.
\]

\item \textbf{Spectral realization by crossing confluence and channel cancellation.}
For any continuous strictly monotone accumulation $E$, if
\[
  E(a)\le c\le E(b),\qquad a\le b,
\]
then there is a unique $t$ such that
\[
  a\le t,\qquad E(t)=c.
\]
At threshold $E(t)=n\pi$,
\[
  \left\lfloor\frac{E(t)}{\pi}\right\rfloor=n.
\]
At a crossing threshold the source event is a confluence: the channel phase
cancels, the amplitude reaches its crest, the resolvent trace consumes the
local residue, the sign changes, and the next harmonic begins on the spectral
readout.  In the source normalization used by the Polya chain,
\[
  E(\Imm\rho)=A_k^\pm,\qquad
  A_0^\pm=\pm\frac{\pi}{2},\qquad
  A_{k+1}^\pm-A_k^\pm=\pm\pi .
\]
So the first crossing on each oriented side costs a half quantum, and every
subsequent threshold is one full $\pi$ step.  The channel cancellation law is:
\[
  \sum_{m\in F} s_m=0
  \Longleftrightarrow
  P_F(s)=M_F(s),
\]
where
\[
  P_F(s)=\sum_{m\in F}\max(s_m,0),\qquad
  M_F(s)=\sum_{m\in F}\max(-s_m,0).
\]
At cancellation,
\[
  \sum_{m\in F}|s_m|=2P_F(s).
\]

\item \textbf{Exhaustion by ladder induction and rigidity.}
If a property holds at $A_0$ and is preserved by $A_k\mapsto A_{k+1}$ and
$A_k\mapsto A_{k-1}$, then it holds at every $A_k$.
Moreover, any crossing sequence $c_n$ satisfying
\[
  E(c_n)=n\pi
\]
is forced to be the purchase-height sequence:
\[
  c_n=\operatorname{purchaseHeight}(n).
\]

\item \textbf{Coordinate transport and Pythagorean readout.}
The rechart
\[
  \operatorname{arcChart}(x)=\frac\pi3x
\]
has inverse
\[
  \operatorname{arcChartInv}(u)=\frac3\pi u.
\]
The critical midpoint is
\[
  \operatorname{arcChart}(x)=\frac\pi6
  \Longleftrightarrow
  x=\frac12.
\]
The value $1/2$ is therefore only the normalized one-dimensional readout of
the source midpoint $\pi/6$.  It is not an independently fitted spectral
constant.  At a source crossing, the standing readout changes sign; that sign
change forces the crossing to occur at the midpoint of the source chart.  The
rule is source-side: projection is only responsible for transporting the event
without losing it or introducing a new one.
The midpoint conclusion routes through the unitary/Pythagorean readout.  In
natural coordinates the readout satisfies
\[
  (\operatorname{readout}_{\sigma,\gamma})_{\Ree}^{2}
  +(\operatorname{readout}_{\sigma,\gamma})_{\Imm}^{2}=1
  \Longleftrightarrow
  \sigma=\frac12.
\]
Equivalently, for the spectral value $w(\rho)=1-1/\rho$,
\[
  \|w(\rho)\|^2=1
  \Longleftrightarrow
  \Ree\rho=\frac12.
\]
The source conservation/no-drift input supplies the unitary source side of this
readout.

\item \textbf{Genuine zero-set object.}
The object is the $L$-function zero set:
\[
  \rho\in\NTZ(\chi)
  \Longleftrightarrow
  0<\Ree\rho<1,\quad L(\rho,\chi)=0.
\]

\item \textbf{Symmetric eigenvalues are real.}
If $T$ is symmetric and $\mu$ is an eigenvalue of $T$, then
\[
  \Imm\mu=0.
\]
The spectral parametrisation is
\[
  \operatorname{spectralZero}(\mu)=\frac12+i\mu .
\]
Therefore a real eigenvalue gives
\[
  \Ree(\operatorname{spectralZero}(\mu))=\frac12 .
\]

\item \textbf{Counting and midpoint identification.}
For the geometric accumulation object,
\[
  \operatorname{harmonicCount}(E,\operatorname{purchaseHeight}(n))=n.
\]
The midpoint chain is
\[
\begin{aligned}
  \operatorname{arcChartInv}(\operatorname{MIDPOINT}_{3D})
    &=\operatorname{MIDPOINT}_{1D}=\frac12,\\
  \operatorname{MIDPOINT}_{3D}
    &=\operatorname{MIDPOINT}_{2D},
\end{aligned}
\]
and the chart/inverse pair is bijective.
\end{enumerate}

\section{Source Crossings and Projection Readout}

\begin{definition}[Source crossing confluence]
For a channel $\chi$ and point $\rho=\sigma+i\gamma$, a source crossing
confluence at height $\gamma$ is a three-dimensional fiber event with the
following simultaneous features:
\begin{enumerate}[label=(C\arabic*)]
\item the channel phases cancel, so the fiber vanishes in the channel sum;
\item the accumulated amplitude reaches the relevant $\pi$-level crest at one
  of the ladder levels $A_m^{\pm}$, with the first oriented crossing carrying
  the half-quantum cost $\pi/2$ and later crossings advancing by full $\pi$
  increments;
\item the crossing consumes the threshold energy required to start the next
  harmonic on the spectral readout, so the outgoing side of the crossing begins
  the next harmonic;
\item the real standing readout changes sign at the crossing;
\item the source midpoint coordinate forced by the sign change and the
  real-axis/unitary readout requirement is $\pi/6$ in the $\pi/3$ chart, while
  the height $\gamma$ is retained as the vertical coordinate;
\item the same crossing location is projected from the three-dimensional helix
  to a two-dimensional unit-circle value $z_2(\rho)$ in the $\pi/3$ chart; the
  subsequent one-dimensional readout is obtained by taking the lifted logarithm
  of that unit-circle value.
\end{enumerate}
Thus the source crossing itself determines the midpoint.  Clause (C5) records
that consequence; it is not an additional lower-dimensional placement
assumption.  After the normalization $x\mapsto 3x/\pi$, this source midpoint is
the value $1/2$.
\end{definition}

Under projection from the three-dimensional fiber to the unit-circle readout,
radial and absolute phase information are collapsed.  The retained datum is the
crossing height $\gamma$, together with the unit-circle event
\[
  z_2(\rho)=\exp(iu_\rho),\qquad u_\rho\in\frac{\pi}{3}\mathbb R .
\]
The angular component of the subsequent one-dimensional readout is the lifted
logarithm
\[
  \Log_{\mathcal H} z_2(\rho)=iu_\rho,
  \qquad
  \ell(\rho)=\frac{3}{\pi}u_\rho
    =
    \frac{3}{\pi}\frac{1}{i}\Log_{\mathcal H}z_2(\rho).
\]
For the midpoint event, $u_\rho=\pi/6$, hence
\[
  \ell(\rho)=\frac{3}{\pi}\cdot\frac{\pi}{6}=\frac12,
\]
so the projected $1/2$ is the logarithmic readout of the $\pi/3$ unit-circle
coordinate.  The branch of $\Log_{\mathcal H}$ is fixed by the continuous helix
lift through the crossing.  The vertical coordinate is recorded separately by
the source-height functional
\[
  \operatorname{height}_{\mathcal H}(z_2(\rho))=\gamma=\Imm\rho,
\]
so the height is retained by the lifted helix branch rather than by the
collapsed radial coordinate.

\begin{definition}[Fiber capture event]
For a Dirichlet channel $\chi$,
\[
  \operatorname{FiberCaptureEvent}_\chi(\rho)
  \quad:\Longleftrightarrow\quad
  \rho\in\NTZ(\chi).
\]
\end{definition}

\begin{definition}[Resolvent pole event]
For a channel $\chi$ and point $\rho$, a resolvent pole event of multiplicity
$m\ge1$ is the local factorization
\[
  L(s,\chi)=(s-\rho)^m g(s),\qquad g(\rho)\ne0,
\]
with $g$ holomorphic near $\rho$.  Equivalently,
\[
  -\frac{L'}{L}(s,\chi)
  =
  -\frac{m}{s-\rho}+H(s),
\]
where $H$ is holomorphic near $\rho$.  In the placed $\pi/3$ trace this is
\[
  \mathcal T_\chi^C(s)
  =
  -\frac{mC^{-\rho}}{s-\rho}
  +H_C(s),
  \qquad C=(\pi/3)^2,
\]
with $H_C$ finite at $\rho$.
\end{definition}

\begin{definition}[Residue accumulator event]
Let $E_\chi$ be the source accumulation ledger for the channel $\chi$.  A
residue accumulator event over $\rho$ consists of:
\begin{enumerate}[label=(A\arabic*)]
\item a resolvent pole event at $\rho$ of multiplicity $m$;
\item a source height $h=\Imm\rho$;
\item a ladder index $k$ and orientation $\varepsilon\in\{+1,-1\}$ such that
\[
  E_\chi(h)=A_k^\varepsilon,\qquad
  A_0^\varepsilon=\varepsilon\frac{\pi}{2},\qquad
  A_{k+1}^\varepsilon-A_k^\varepsilon=\varepsilon\pi;
\]
\item consumption of the trace residue at that threshold, i.e. the pole
residue is assigned to the threshold crossing at height $h$.
\end{enumerate}
Thus the residue accumulator is not an additional zero-locating formula.  It
is the source ledger that records where the already identified trace residue is
spent in the helix.
\end{definition}

\begin{definition}[Spectral harmonic event]
A spectral harmonic event over $\rho$ is a residue accumulator event whose
threshold starts the next harmonic on the spectral readout.  Equivalently, if
$h=\Imm\rho$ is the threshold height and $E_\chi(h)=n\pi$ in the positive
orientation, then
\[
  \operatorname{harmonicCount}(E_\chi,h)=n.
\]
For the negative orientation the same statement holds after applying the
reflected anti-helix orientation.
\end{definition}

\begin{definition}[Projection readout event]
A projection readout event over $\rho$ is a source crossing confluence whose
unit-circle projection is
\[
  z_2(\rho)=\exp(iu_\rho)
\]
on the continuous helix branch, with $u_\rho=\pi/6$ at the crossing midpoint.
Its one-dimensional readout is
\[
  \ell(\rho)=\frac3\pi u_\rho=\frac12,
\]
and its height is retained by the lifted phase:
\[
  \operatorname{height}_{\mathcal H}(z_2(\rho))=\Imm\rho .
\]
\end{definition}

\begin{definition}[Source fiber event]
For a non-principal Dirichlet channel $\chi$ and point $\rho$, a source fiber
event is a three-dimensional event over $\rho$ carrying the following data:
\begin{enumerate}[label=(S\arabic*)]
\item a source drift rate $\lambda_\rho$ and nonzero amplitude $a_\rho$ with
\[
  e^{\Ree(\lambda_\rho)\tau}a_\rho=a_\rho
  \qquad(\forall \tau\in\mathbb R);
\]
\item the $\pi/3$ source coordinate
\[
  x_\rho=\operatorname{arcChart}(\Ree\rho)=\frac{\pi}{3}\Ree\rho;
\]
\item height retention
\[
  \Imm(\lambda_\rho)=\Imm\rho;
\]
\item nonzero spectral value $\rho\ne0$;
\item source-chain cancellation with explicit decay;
\item the resolvent-trace pole at $\rho$.
\end{enumerate}
It does not contain the conclusion $\Ree\rho=1/2$.
\end{definition}

\begin{definition}[Source-resolved zero]
A source-resolved zero for a channel $\chi$ is a point $\rho$ carrying all of
the following data:
\[
\begin{gathered}
  \operatorname{FiberCaptureEvent}_\chi(\rho),\qquad
  \operatorname{SourceFiberEvent}_\chi(\rho),\\
  \operatorname{ResolventPoleEvent}_\chi(\rho),\qquad
  \operatorname{ResidueAccumulatorEvent}_\chi(\rho),\\
  \operatorname{SpectralHarmonicEvent}_\chi(\rho),\qquad
  \operatorname{ProjectionReadoutEvent}_\chi(\rho).
\end{gathered}
\]
This is the precise object closed by the canonical proof: it is an analytic
zero together with its source trace, source crossing, spectral harmonic, and
projection readout.  The one-dimensional point produced by that readout is
identified with the same zero $\rho$:
\[
  \rho=\ell(\rho)+i\,\operatorname{height}_{\mathcal H}(z_2(\rho)).
\]
\end{definition}

\section{Canonical Proof Spine}

This section states the proof directly.  The objects
are the prime-generated source helix, its channel fibers, the resolvent trace,
and the spectral operator induced by the Gram form.

\begin{definition}[Canonical $\pi/3$ chart]
Let
\[
  \alpha(x)=\frac{\pi}{3}x,\qquad
  \alpha^{-1}(u)=\frac3\pi u.
\]
The source and unit-circle midpoint is
\[
  m_{3D}=m_{2D}=\frac{\pi}{6},
\]
and the normalized one-dimensional midpoint is
\[
  m_{1D}=\alpha^{-1}(m_{3D})=\frac12 .
\]
Thus
\[
  \alpha(x)=\frac{\pi}{6}
  \Longleftrightarrow
  x=\frac12 .
\]
\end{definition}

Throughout the proof, a source crossing means the confluence event defined in
the previous section: phase cancellation, amplitude threshold, sign change,
new harmonic creation, and projection by lifted phase.  Its first oriented
cost is $\pi/2$, and later crossings advance by full $\pi$ increments:
\[
  A_0^\pm=\pm\frac{\pi}{2},\qquad
  A_{k+1}^\pm-A_k^\pm=\pm\pi .
\]

\begin{lemma}[No radial drift]
The source helix has no radial drift.  If a nonzero source amplitude $c$ obeys
\[
  e^{\Ree\lambda\,\tau}c=c
  \qquad(\forall \tau\in\mathbb R),
\]
then
\[
  \Ree\lambda=0.
\]
\end{lemma}

\begin{proof}
Since $c\ne0$, division by $c$ gives
\[
  e^{\Ree\lambda\,\tau}=1
  \qquad(\forall \tau).
\]
Taking logarithmic derivative in $\tau$ gives $\Ree\lambda=0$.  Equivalently,
the only exponential one-parameter scaling that fixes a nonzero amplitude for
all time is the trivial scaling.
\end{proof}

\begin{lemma}[Threshold rigidity]
Let $E:\mathbb R\to\mathbb R$ be the monotone accumulation function of the
source helix.  For each $n\ge0$ there is a unique purchase height $h_n$ with
\[
  E(h_n)=n\pi.
\]
Moreover any nonnegative sequence $c_n$ satisfying
\[
  E(c_n)=n\pi
\]
coincides with $h_n$ for all $n$.
\end{lemma}

\begin{proof}
Existence and uniqueness follow from continuity and strict monotonicity of
$E$.  Rigidity follows because strict monotonicity makes every level set of
$E$ a singleton.
\end{proof}

\begin{lemma}[Harmonic count at threshold]
At the $n$-th purchase height,
\[
  \operatorname{harmonicCount}(E,h_n)=n.
\]
\end{lemma}

\begin{proof}
The harmonic count is the threshold counter for the ladder $E(t)=n\pi$.
By threshold rigidity, the $n$-th threshold is reached exactly at $h_n$, so the
counter reads $n$ there.
\end{proof}

\begin{lemma}[Crossing forces the midpoint]
Every source crossing occurs at the source midpoint
\[
  m_{3D}=\frac{\pi}{6}.
\]
After projection to the unit-circle chart and logarithmic readout,
\[
  m_{3D}=m_{2D},\qquad
  \alpha^{-1}(m_{2D})=\frac12 .
\]
\end{lemma}

\begin{proof}
At a crossing, the standing readout changes sign.  The source readout is real
at the crossing and has no radial drift by the previous lemma, so the sign
change is an event on the real midpoint axis of the $\pi/3$ chart.  The midpoint
of the source chart is $\pi/6$.  The projection to the unit-circle chart
preserves the angular midpoint, and applying $\alpha^{-1}$ gives
\[
  \alpha^{-1}\!\left(\frac{\pi}{6}\right)
  =
  \frac3\pi\cdot\frac{\pi}{6}
  =
  \frac12 .
\]
\end{proof}

\begin{lemma}[Resolvent pole at a captured zero]
Let $\rho\in\NTZ(\chi)$.  Then the logarithmic derivative trace
\[
  -\frac{L'}{L}(s,\chi)
\]
has no finite punctured limit at $s=\rho$.  More precisely, if
\[
  m=\operatorname{ord}_\rho L(\cdot,\chi)\ge1,
\]
then near $\rho$
\[
  \frac{L'}{L}(s,\chi)
  =
  \frac{m}{s-\rho}+H(s),
\]
where $H$ is holomorphic near $\rho$.
\end{lemma}

\begin{proof}
Write
\[
  L(s,\chi)=(s-\rho)^m g(s),
\]
where $g$ is holomorphic near $\rho$ and $g(\rho)\ne0$.  Then
\[
  \frac{L'}{L}(s,\chi)
  =
  \frac{m}{s-\rho}+\frac{g'}{g}(s),
\]
and $g'/g$ is holomorphic near $\rho$.  Thus $-L'/L$ has a pole at $\rho$ with
residue $-m$ and no finite punctured limit.
\end{proof}

\begin{lemma}[Prime trace identity]
In the half-plane of absolute convergence,
\[
  \sum_{n\ge1}\chi(n)\Lambda(n)n^{-s}
  =
  -\frac{L'}{L}(s,\chi).
\]
With the $\pi/3$ gauge $C=(\pi/3)^2$, the placed trace is
\[
  \mathcal T_\chi^C(s)
  =
  C^{-s}\left(-\frac{L'}{L}(s,\chi)\right).
\]
\end{lemma}

\begin{proof}
Differentiate the Euler product logarithm:
\[
  L(s,\chi)=\prod_p(1-\chi(p)p^{-s})^{-1}.
\]
For $\Ree s>1$,
\[
  \log L(s,\chi)
  =
  \sum_p\sum_{k\ge1}\frac{\chi(p)^k}{k}p^{-ks}.
\]
Differentiation gives
\[
  -\frac{L'}{L}(s,\chi)
  =
  \sum_p\sum_{k\ge1}\chi(p)^k(\log p)p^{-ks}
  =
  \sum_{n\ge1}\chi(n)\Lambda(n)n^{-s}.
\]
Multiplication by $C^{-s}$ is exactly the placed $\pi/3$ source gauge.
\end{proof}

\begin{lemma}[Captured zero gives a resolvent pole event]
If $\rho\in\NTZ(\chi)$ has vanishing order $m\ge1$, then $\rho$ determines a
resolvent pole event of multiplicity $m$.
\end{lemma}

\begin{proof}
This is exactly the local factorization in the resolvent pole lemma.  The
placed trace formula follows by multiplying the principal part by the nonzero
gauge $C^{-s}$ and evaluating the residue factor at $s=\rho$.
\end{proof}

\begin{lemma}[Prime trace places the pole on the source]
In the half-plane where the Euler product converges, the pole readout is the
continuation of the prime-power source trace
\[
  \mathcal T_\chi^C(s)
  =
  \sum_{n\ge1}\Lambda(n)\chi(n)(Cn)^{-s}.
\]
Thus a resolvent pole event is a source-trace event once the source fiber is
identified by the Euler-product discriminator.
\end{lemma}

\begin{proof}
The prime trace identity gives
\[
  \mathcal T_\chi^C(s)=C^{-s}\left(-\frac{L'}{L}(s,\chi)\right).
\]
The Euler-product discriminator supplies the arithmetic source of the trace:
the coefficients are the prime-power von Mangoldt weights generated by the
Euler product.  The gauge $C^{-s}$ is nonzero, so it transports the pole and
its multiplicity without changing the event location.
\end{proof}

\begin{lemma}[Accumulator threshold produces the crossing and harmonic]
Let $\rho$ carry a residue accumulator event.  Then the source height
$h=\Imm\rho$ is a threshold crossing.  If $E_\chi(h)=n\pi$ in the positive
orientation, then the corresponding spectral readout begins the $n$-th
harmonic:
\[
  \operatorname{harmonicCount}(E_\chi,h)=n.
\]
\end{lemma}

\begin{proof}
By definition of the residue accumulator event, the trace residue is consumed
at a ladder threshold.  Threshold rigidity makes the threshold height the
unique purchase height for its level.  The harmonic-count lemma then gives the
displayed harmonic count.  The source crossing confluence is precisely the
same threshold read with phase cancellation, amplitude cresting, sign change,
and outgoing harmonic creation.
\end{proof}

\begin{lemma}[Crossing projection gives the midpoint readout]
Let $\rho$ carry a source crossing confluence and projection readout event.
Then the projected one-dimensional real coordinate is $1/2$ and the height is
retained by the lifted phase.
\end{lemma}

\begin{proof}
The crossing forces the source midpoint $\pi/6$.  Projection preserves this
midpoint in the unit-circle chart.  Applying the inverse $\pi/3$ chart gives
\[
  \alpha^{-1}\!\left(\frac{\pi}{6}\right)
  =
  \frac3\pi\cdot\frac{\pi}{6}
  =
  \frac12 .
\]
The lifted branch $\Log_{\mathcal H}$ records the vertical phase through the
crossing, and the source-height functional records
\[
  \operatorname{height}_{\mathcal H}(z_2(\rho))=\Imm\rho .
\]
\end{proof}

\begin{proposition}[Correspondence chain for a source-resolved zero]
For a source-resolved zero $\rho$, the zero-side pole, source crossing, and
spectral harmonic are the same event read through three ledgers:
\[
  \text{resolvent trace}
  \quad\longrightarrow\quad
  \text{source accumulator}
  \quad\longrightarrow\quad
  \text{spectral harmonic}.
\]
\end{proposition}

\begin{proof}
The source-resolved zero includes a fiber capture event, so the captured-zero
lemma gives the resolvent pole.  The prime-trace lemma places this pole on the
prime-generated source trace.  The residue accumulator lemma identifies the
source threshold and starts the harmonic.  The projection lemma then gives the
midpoint readout and retained height.  These are not separate zeros; they are
the same source event expressed in the trace, source, spectral, and projection
coordinates.
\end{proof}

\begin{theorem}[Canonical closure for source-resolved zeros]
Every source-resolved zero for the prime-generated source helix has
one-dimensional real readout $1/2$:
\[
  \operatorname{SourceResolvedZero}_\chi(\rho)
  \Longrightarrow
  \Ree\rho=\frac12 .
\]
\end{theorem}

\begin{proof}
Let $\rho$ be a source-resolved zero.  By the correspondence chain, $\rho$
carries the source crossing and projection readout events.  The crossing
projection lemma gives the normalized midpoint readout
\[
  \frac3\pi\cdot\frac{\pi}{6}=\frac12 .
\]
Since the source-resolved zero identifies the projected one-dimensional readout
with the real coordinate of $\rho$, it follows that $\Ree\rho=1/2$.
\end{proof}

\begin{corollary}[Zeta/L1 source-resolved channel]
For the principal zeta channel, use the eta-regularized source fiber
\[
  F_{L1}(s)=C^{-s}(1-2^{1-s})\zeta(s).
\]
Off $\Ree s=1$ the eta factor is nonzero, so the L1 fiber zeros are exactly the
zeta zeros.  Hence every source-resolved nontrivial zeta zero projects to real
coordinate $1/2$.
\end{corollary}

\begin{proof}
The eta factor vanishes only when $2^{1-s}=1$, which forces $\Ree s=1$.  In the
critical strip away from $\Ree s=1$, it is nonzero.  Thus the L1 fiber has the
same nontrivial zeros as $\zeta$, and the canonical closure applies to the
eta-regularized principal source.
\end{proof}

\section{Mathematical Reproduction of the Computation}

The computation evaluates the source fiber equations above.  It is not a proof
step.  The crossing ordinates are defined by the complete geometric fiber and
its standing readout.  Finite summation, residue decomposition, Hurwitz
notation, and Euler--Maclaurin summation appear only where the derived fiber is
represented for one-dimensional evaluation \cite{Apostol1976}; they are not
sources of the fiber.

\subsection{Placed Geometric Fiber}

The geometry is
\[
  A=\pi/6,\qquad ds=\pi/3,\qquad C=2A\,ds=(\pi/3)^2.
\]
For $s=\frac12+it$, the placed finite head is
\[
  H_M(t)
  =
  \sum_{n=1}^M \chi(n)R_n^{-1}e^{-it\Log(R_n^2)},
  \qquad R_n^2=Cn.
\]
Thus
\[
  H_M(t)=\sum_{n=1}^M \chi(n)(Cn)^{-s}.
\]
The tail is not added from a separate external object.  It is the remaining
geometric fiber beyond the cutoff:
\[
  G_M(t)
  =
  \sum_{n>M}\chi(n)R_n^{-1}e^{-it\Log(R_n^2)}
  =
  \sum_{n>M}\chi(n)(Cn)^{-s}.
\]
The complete placed fiber is
\[
  F_\chi(t)=H_M(t)+G_M(t),
\]
independently of the cutoff $M$ when the tail is included.

\subsection{Residue Coordinates for the Geometric Tail}

The geometric tail may be written in conductor residue coordinates.  Let
$n_{0,r}$ be the first integer $>M$ in residue class $r\bmod q$.  Then
\[
  G_M(t)
  =
  \sum_r \chi(r) C^{-s}q^{-s}
  \sum_{m\ge0}\left(\frac{n_{0,r}}q+m\right)^{-s}.
\]
This identity is just the residue-class reindexing of the geometric terms
$R_n^2=Cn$.  Hurwitz notation abbreviates the inner residue tail; it does not
define the tail or the crossing condition.  It is an auxiliary representation
of the derived fiber.

For finite one-dimensional evaluation, replace the inner infinite residue tail by
\[
  \mathcal T_{h,J}(s,x)
  \approx
  \sum_{m\ge0}(x+m)^{-s}.
\]
For integers $h,J\ge1$, set $X=x+h$ and write
\[
  (s)^{\overline m}=s(s+1)\cdots(s+m-1),\qquad
  (s)^{\overline0}=1 .
\]
The Euler--Maclaurin summation evaluator used for that geometric residue tail is
\[
\begin{aligned}
  \mathcal T_{h,J}(s,x)
  &=
  \sum_{m=0}^{h-1}(x+m)^{-s}
  +\frac{X^{1-s}}{s-1}
  +\frac12 X^{-s}  \\
  &\quad+
  \sum_{j=1}^{J}
    \frac{B_{2j}}{(2j)!}
    (s)^{\overline{2j-1}}
    X^{-s-2j+1},
\end{aligned}
\]
where $B_{2j}$ are Bernoulli numbers.  The associated remainder is the standard
Euler--Maclaurin remainder after the $B_{2J}$ term.  The numerical evaluations
use $h=8$ and finite $J$ chosen large enough for the working
precision.

With this evaluator, the finite evaluated fiber is
\[
  F_{M,\chi}(t)
  =
  \sum_{n=1}^M \chi(n)(Cn)^{-s}
  +
  \sum_r \chi(r) C^{-s}q^{-s}
  \mathcal T_{h,J}\!\left(s,\frac{n_{0,r}}q\right).
\]

For the L1 channel, the same geometric tail is grouped with conductor $2$ and
weights
\[
  \chi_\eta(1)=1,\qquad \chi_\eta(0)=-1,
\]
so that
\[
  F_{M,L1}(t)
  =
  C^{-s}2^{-s}
  \left[
    \mathcal T_{h,J}\!\left(s,\frac{n_{0,1}}2\right)
    -
    \mathcal T_{h,J}\!\left(s,\frac{n_{0,0}}2\right)
  \right]
  +
  \sum_{n=1}^{M}(-1)^{n+1}(Cn)^{-s}.
\]
This is the finite evaluation of the eta-regularized zeta fiber in placed
geometric coordinates.

\subsection{Principal One-Dimensional Tail Error Control}

For the projected principal/L1 tail, use the classical Euler--Maclaurin
summation identity as an error-control theorem \cite{WhittakerWatson1927}.
Let
\[
  S_N(s)=\sum_{n=1}^{N}n^{-s},\qquad \{u\}=u-\lfloor u\rfloor ,
\]
and define
\[
  I(s,N)=\int_N^\infty \{u\}\,u^{-s-1}\,du,\qquad
  R(s,N)=sI(s,N).
\]
For $\Ree s>0$ and $s\ne1$, the continued one-dimensional response is
\[
  c(s)
  =
  \frac{s}{s-1}
  -
  s\int_1^\infty \{u\}\,u^{-s-1}\,du,
\]
and the Euler--Maclaurin identity is
\[
  c(s)=\zeta(s),
\]
with finite-readout error
\[
  S_N(s)-\zeta(s)-\frac{N^{1-s}}{1-s}=R(s,N)
\]
and bound
\[
  \|R(s,N)\|
  \le
  \frac{\|s\|}{\Ree s}\,N^{-\Ree s}
  \qquad (N\ge2).
\]
This theorem controls the finite projected evaluation of the principal
geometric tail; it is a representation of the derived fiber, not a source for
the tail.

The full finite readout is
\[
  Z_\chi(t)=\Ree(e^{i\Theta_\chi(t)}F_{M,\chi}(t)).
\]
Crossings are isolated by sign changes
\[
  Z_\chi(a)Z_\chi(b)<0
\]
followed by bisection.  If $\gamma\in(a,b)$ is a simple crossing,
\[
  Z_\chi(\gamma)=0,\qquad Z_\chi'(\gamma)\ne0,
\]
then there is an interval around $\gamma$ on which the sign changes once, and
bisection on any bracket inside that interval converges to $\gamma$.

For local reproduction from an intermediate height $t_0$, choose $\delta>0$ and
evaluate $Z_\chi$ on $[t_0-\delta,t_0+\delta]$.  If a sign change occurs in
$[a,b]\subset[t_0-\delta,t_0+\delta]$, bisection on $Z_\chi$ rederives the
crossing in that local interval; no earlier crossing is used.

\bibliographystyle{alpha}
\bibliography{proof_bridge_refs}

\end{document}
